\section{Serialize and Deserialize a Binary Tree}
\label{sec:trees-serialize}

\problemstatement{Design two functions: \emph{serialize}, which encodes a
binary tree into a single string, and \emph{deserialize}, which rebuilds
the exact tree from that string. Values may repeat and the tree may be
arbitrary.}

\begin{patternbox}
Unlike reconstruction from two traversals, here we control the format --
so we choose one that is \emph{self-describing} by \textbf{recording null
children explicitly}. With null markers, a \emph{single} preorder walk
carries all the structure needed to rebuild the tree unambiguously. The
design freedom is the whole point.
\end{patternbox}

\subsection*{Understanding the problem carefully}

Two traversals are needed to rebuild a tree only because a normal traversal
hides where subtrees end. If we instead emit a special token (say \verb|#|)
whenever we hit a null child, the string records the tree's exact shape.
Reading that string back in the same preorder order, a null token means
``this child is empty,'' so deserialization mirrors serialization step for
step -- read a token, build a node (or return null), then recursively build
its left and right children.

\begin{ideabox}[Preorder with Explicit Null Markers Is Self-Sufficient]
\textbf{Serialize:} preorder traversal; output each value, and output a
sentinel \verb|#| for every null child, separated by commas.
\textbf{Deserialize:} read tokens in order; \verb|#| yields null,
otherwise create a node and recursively deserialize its left then right
child. The null markers tell the reader exactly where each subtree stops.
\end{ideabox}

\subsection*{Algorithm}

\begin{enumerate}
  \item \textsc{Serialize}(node, out): if node is null, append
        \verb|#,|; else append \verb|node.val,| then
        \textsc{Serialize}(left) and \textsc{Serialize}(right).
  \item \textsc{Deserialize}: read the next token; if \verb|#| return
        null; else make a node, set node.left $=$ \textsc{Deserialize},
        node.right $=$ \textsc{Deserialize}; return node.
\end{enumerate}

\subsection*{Dry run}

\dryrun{Tree: root $1$; left $2$; right $3$ (with children $4,5$).}

\begin{itemize}
  \item Serialize: \verb|1,2,#,#,3,4,#,#,5,#,#,| -- $2$ is a leaf (two
        \verb|#|), then $3$, then its children $4,5$.
  \item Deserialize reads $1$ (node), builds left from $2$ (then two
        \verb|#| $\to$ leaf), builds right from $3$, whose children come
        from $4$ and $5$.
  \item The rebuilt tree is identical to the original.
\end{itemize}

\subsection*{C++ implementation}

\begin{lstlisting}
#include <bits/stdc++.h>
using namespace std;

struct TreeNode {
    int val;
    TreeNode* left;
    TreeNode* right;
    TreeNode(int v) : val(v), left(nullptr), right(nullptr) {}
};

void serialize(TreeNode* node, string& out) {
    if (node == nullptr) { out += "#,"; return; }
    out += to_string(node->val) + ",";
    serialize(node->left, out);
    serialize(node->right, out);
}

TreeNode* deserialize(istringstream& in) {
    string token;
    getline(in, token, ',');
    if (token == "#") return nullptr;
    TreeNode* node = new TreeNode(stoi(token));
    node->left  = deserialize(in);
    node->right = deserialize(in);
    return node;
}

void inorderPrint(TreeNode* n) {
    if (!n) return;
    inorderPrint(n->left);
    cout << n->val << " ";
    inorderPrint(n->right);
}

int main() {
    TreeNode* root = new TreeNode(1);
    root->left = new TreeNode(2);
    root->right = new TreeNode(3);
    root->right->left = new TreeNode(4);
    root->right->right = new TreeNode(5);

    string s;
    serialize(root, s);
    cout << "Serialized: " << s << "\n";

    istringstream in(s);
    TreeNode* rebuilt = deserialize(in);
    cout << "Inorder of rebuilt: ";
    inorderPrint(rebuilt);
    cout << "\n";
    return 0;
}
\end{lstlisting}

\begin{complexitybox}
\textbf{Time Complexity.} $O(n)$ for both serialize and deserialize --
each node and each null marker is written and read once. \textbf{Space
Complexity.} $O(n)$ for the string, plus $O(h)$ recursion.
\end{complexitybox}

\begin{takeawaybox}
When you own the encoding, adding explicit null markers makes a single
traversal invertible -- no second traversal required. This ``sentinel for
absence'' idea is the general trick for making any recursive structure
round-trip through a flat string.
\end{takeawaybox}
